\section{Tensor conductor energy with the mask retained}
\label{sec:tensor-energy}

For \(u=ga\), \(v=gb\), set
\begin{equation}
 C_{g,a,b}(\alpha,\beta)
 :=\sum_{m,n}c_{m,n}^{(g)}\alpha(m)\beta(n),
 \qquad
 \alpha\in\mathfrak X(u),\quad\beta\in\mathfrak X(b).
 \label{eq:rectangular-C}
\end{equation}
By \cref{thm:pairwise-rectangularization}, this is the full masked pair
transform after removal of the harmless phase
\(\overline{\chi(-h)}\).

Put
\begin{equation}
 \cK_Q(F):=F+\frac QF.
 \label{eq:KQ}
\end{equation}
For dyadic coordinate-conductor sizes \(F_1,F_2\), define the
pair-separated energy
\begin{align}
 \cE_Y(F_1,F_2)
 :={}&
 \sum_{\substack{g,a,b\ {\rm pairwise\ coprime}\\
                  ga,gb\le R,\ gab\asymp Y\\
                  gab\ {\rm squarefree},\ (gab,H)=1}}
 \frac{|\vartheta_R(ga)\vartheta_R(gb)|^2}{\varphi(gab)}
 \notag\\
 &\times
 \sum_{\substack{\alpha\in\mathfrak X(ga)\\\cond\alpha\asymp F_1}}
 \sum_{\substack{\beta\in\mathfrak X(b)\\\cond\beta\asymp F_2}}
 |C_{g,a,b}(\alpha,\beta)|^2.
 \label{eq:pair-energy}
\end{align}
This is intentionally a sum of pair energies, not the energy of a
fictional global rectangular projection of \(Z_q\).

\begin{theorem}[Arbitrary-matrix tensor conductor bound]
\label{thm:tensor-energy}
Uniformly for \(Y\le R^2\) and \(1\le F_i\le R\),
\begin{equation}
 \boxed{
 \cE_Y(F_1,F_2)
 \ll_\eps X^\eps Q^2
 \cK_Q(F_1)\cK_Q(F_2).}
 \label{eq:tensor-energy-bound}
\end{equation}
No separability assumption on \(c_{m,n}^{(g)}\) is made.
\end{theorem}

\begin{proof}
Use \eqref{eq:theta-bound}, split \(u=ga\) and \(b\) into dyadic
ranges, and fix \(g\).  Since \(q=gab=ub\) and \((u,b)=1\),
\(\varphi(q)=\varphi(u)\varphi(b)\).  Retaining squarefreeness but
removing the conditions
\[
 g\mid u,\qquad (u,b)=1,\qquad ub\asymp Y,\qquad gb\le R
\]
enlarges the nonnegative sum.  It is therefore enough to bound
\begin{align}
 &\sum_{\substack{u\asymp U\\u\ {\rm squarefree}}}\frac1{\varphi(u)}
 \sum_{\substack{\alpha\in\mathfrak X(u)\\\cond\alpha\asymp F_1}}
 \sum_{\substack{b\asymp B\\b\ {\rm squarefree}}}\frac1{\varphi(b)}
 \sum_{\substack{\beta\in\mathfrak X(b)\\\cond\beta\asymp F_2}}
 \left|\sum_{m,n}c_{m,n}^{(g)}\alpha(m)\beta(n)\right|^2.
 \label{eq:enlarged-tensor-sum}
\end{align}

Fix \(b,\beta\) and apply
\cref{thm:exact-conductor-large-sieve} in \(u,\alpha\) to
\[
 a_m(\beta)=\sum_n c_{m,n}^{(g)}\beta(n).
\]
This contributes \(\cK_Q(F_1)\) and leaves
\[
 \cK_Q(F_1)
 \sum_{\substack{b\asymp B\\b\ {\rm squarefree}}}\frac1{\varphi(b)}
 \sum_{\cond\beta\asymp F_2}
 \sum_m\left|\sum_nc_{m,n}^{(g)}\beta(n)\right|^2.
\]
Apply the same theorem in \(b,\beta\), separately for each \(m\).
The result is
\begin{equation}
 \ll_\eps X^\eps\cK_Q(F_1)\cK_Q(F_2)
 \sum_{m,n}|c_{m,n}^{(g)}|^2.
 \label{eq:fixed-g-tensor-bound}
\end{equation}
This is a composition of two Hilbert-space operator bounds, so it does
not factor the coefficient matrix.

Sum \eqref{eq:fixed-g-tensor-bound} over \(g\) and the logarithmically
many dyadic \(u,b\)-ranges.  Then
\cref{prop:compatibility-energy} proves
\eqref{eq:tensor-energy-bound}.
\end{proof}

\subsection{The inherited norm route}

For a fixed pair, multiplicative Parseval and
\cref{cor:product-conductors} decompose its centered fiber into the
blocks \((F_1,F_2)\ne(1,1)\).  Let
\(D_{u,v}^{F_1,F_2}\) denote the inverse character transform of one such
block, with the factor
\(\vartheta_R(u)\vartheta_R(v)\) included.  Then
\begin{equation}
 \|D_{u,v}^{F_1,F_2}\|_{2,\cU_q}^2
 =\frac{|\vartheta_R(u)\vartheta_R(v)|^2}{\varphi(q)}
 \sum_{\substack{\cond\alpha\asymp F_1\\
                  \cond\beta\asymp F_2}}
 |C_{g,a,b}(\alpha,\beta)|^2.
 \label{eq:pair-block-parseval}
\end{equation}
The phase \(\overline{\chi(-h)}\) has modulus one.

Let \(\cL_{u,v}^{F_1,F_2}\) be the scalar packet
\eqref{eq:scalar-packet} formed from this pair block, and define the
pair-separated dyadic aggregate
\begin{equation}
 \mathfrak L_Y(F_1,F_2)
 :=\sum_{\substack{g,a,b\ {\rm pairwise\ coprime}\\
                   ga,gb\le R,\ gab\asymp Y\\
                   gab\ {\rm squarefree},\ (gab,H)=1}}
 |\cL_{ga,gb}^{F_1,F_2}|.
 \label{eq:pair-separated-scalar-block}
\end{equation}
This is one object.  The subscripts \(A\) and \(B\) used informally
below name two estimates for it, not two different projections.  By
linearity, pairwise triangle inequality, and dyadic conductor
decomposition,
\begin{equation}
 \sum_{q\asymp Y}|\cL_q^{\rm disc}|
 \le\sum_{\substack{F_1,F_2\\(F_1,F_2)\ne(1,1)}}
 \mathfrak L_Y(F_1,F_2).
 \label{eq:block-reassembly-inequality}
\end{equation}

The number of triples \((g,a,b)\) with \(gab\asymp Y\) is
\(\ll YX^\eps\), by \eqref{eq:lcm-pair-count}.  Cauchy--Schwarz and
\cref{thm:tensor-energy} therefore give
\begin{equation}
 \sum_{\substack{q\asymp Y\\\operatorname{lcm}(u,v)=q}}
 \|D_{u,v}^{F_1,F_2}\|_{2,\cU_q}
 \ll_\eps
 X^\eps Q\sqrt{Y\cK_Q(F_1)\cK_Q(F_2)}.
 \label{eq:block-fiber-norm}
\end{equation}
The notation on the left means the sum over the pair contributions
before their lcm reassembly.

\begin{corollary}[Tensor-norm scalar bound]
\label{cor:tensor-norm-route}
Under the inherited hypotheses of \eqref{eq:inherited-norm-route}, every
nonprincipal coordinate-conductor block satisfies
\begin{equation}
 \boxed{
 \mathfrak L_Y(F_1,F_2)
 \ll_\eps X^\eps Q
 \sqrt{JY\cK_Q(F_1)\cK_Q(F_2)}.}
 \label{eq:tensor-norm-scalar}
\end{equation}
\end{corollary}

\begin{proof}
Use \eqref{eq:centered-pair-sum} and Minkowski before
\eqref{eq:inherited-norm-route}, then insert
\eqref{eq:block-fiber-norm}.  The Fourier-tail error is already part of
the inherited polynomial-mass hypothesis.
\end{proof}

The pairwise triangle inequality may discard cancellation between lcm
decompositions, but it costs no fixed power and is sufficient for the
claimed closure.  The next section proves a second estimate for this
same aggregate.
